% from ega0-16-fr.tex
\oldpage[0\textsubscript{IV-1}]{36}
\begin{enumerate}
  \item[(i)] Pour tout élément $M$-régulier $x$ appartenant à l'idéal maximal
  de $A$, on a
  $\operatorname{coprof}(M/xM)=\operatorname{coprof}(M)$.
  \item[(ii)] On a
  $\operatorname{coprof}_A(M)=\operatorname{coprof}_{\widehat A}(\widehat M)$.
\end{enumerate}
\end{proposition}

\begin{proof}
En effet, (i) résulte de (16.3.4) et (16.4.6, (ii)), et (ii) de (16.2.4) et
(16.4.6, (iv)).
\end{proof}

Nous montrerons plus tard (IV, 6.11.5) que pour tout idéal premier
$\mathfrak{p}$ de $A$, on a
$\operatorname{coprof}_{A_{\mathfrak{p}}}(M_{\mathfrak{p}})
\leq\operatorname{coprof}_A(M)$.

\begin{proposition}[16.4.11]
\label{0.16.4.11-fr}
Sous les hypothèses de (16.4.8), on a
\begin{equation}
\label{0.16.4.11.1-fr}
\tag{16.4.11.1}
  \operatorname{coprof}_A(M_{[\rho]})=\operatorname{coprof}_B(M).
\end{equation}
\end{proposition}

\begin{proof}
Cela résulte de (16.1.9) et (16.4.8).
\end{proof}

\subsection{Modules de Cohen-Macaulay}
\label{subsection:0.16.5-fr}

\begin{definition}[16.5.1]
\label{0.16.5.1-fr}
Soit $A$ un anneau local noethérien. On dit qu'un $A$-module de type fini $M$
est un \emph{module de Cohen-Macaulay} si
$\operatorname{coprof}(M)=0$, autrement dit si $M=0$ ou si $M\neq0$ et
$\dim(M)=\operatorname{prof}(M)$. On dit que $A$ est un
\emph{anneau de Cohen-Macaulay} si le $A$-module $A$ est un module de
Cohen-Macaulay.
\end{definition}

Un $A$-module de longueur finie est un module de Cohen-Macaulay puisqu'il est
de dimension $0$ (16.1.10). Un anneau local réduit $A$ de dimension $1$ est un
anneau de Cohen-Macaulay, car alors son idéal maximal $\mathfrak{m}$ ne peut
être associé à $A$ (Bourbaki, \emph{Alg.~comm.}, chap.~IV, \S~2, no~5,
prop.~10), donc $\operatorname{prof}(A)=1$ en vertu de (16.4.6, (i)). Si $A$
est un \emph{anneau local intègre et intégralement clos de dimension $\geq2$},
on a $\operatorname{prof}(A)\geq2$~: en effet, si $x\neq0$ est un élément de
$\mathfrak{m}$, on sait (Bourbaki, \emph{Alg.~comm.}, chap.~VII, \S~1, no~4,
prop.~8) que les idéaux premiers associés à $A/xA$ sont de hauteur $1$, donc
distincts de $\mathfrak{m}$, et leur réunion est par suite distincte de
$\mathfrak{m}$ (Bourbaki, \emph{Alg.~comm.}, chap.~II, \S~1, no~1, prop.~2)~;
il existe donc des suites $(x,y)$ qui sont $A$-régulières et formées
d'éléments de $\mathfrak{m}$, d'où notre assertion. On en conclut qu'un anneau
local intégralement clos de dimension $2$ est un anneau de Cohen-Macaulay.

\begin{proposition}[16.5.2]
\label{0.16.5.2-fr}
Soient $A$ un anneau local noethérien, $M$ un $A$-module de type fini. Pour que
$M$ soit un module de Cohen-Macaulay, il faut et il suffit que $\widehat M$
soit un $\widehat A$-module de Cohen-Macaulay.
\end{proposition}

\begin{proof}
Cela résulte aussitôt de (16.4.10, (ii)).
\end{proof}

\begin{proposition}[16.5.3]
\label{0.16.5.3-fr}
Sous les hypothèses de (16.4.8), pour que $M$ soit un $B$-module de
Cohen-Macaulay, il faut et il suffit que $M_{[\rho]}$ soit un $A$-module de
Cohen-Macaulay.
\end{proposition}

\begin{proof}
Cela résulte de (16.4.11).
\end{proof}

\begin{proposition}[16.5.4]
\label{0.16.5.4-fr}
Soient $A$ un anneau local noethérien, $M$ un $A$-module de type fini $\neq0$.
Si $M$ est un $A$-module de Cohen-Macaulay, on a
$\dim(A/\mathfrak{p})=\dim(M)$ pour tout idéal premier
$\mathfrak{p}\in\operatorname{Ass}(M)$~; en particulier aucun idéal premier
associé à $M$ n'est immergé.
\end{proposition}

\begin{proof}
Cela résulte aussitôt des inégalités (16.4.6.2).
\end{proof}

On peut encore dire que $M$ (ou $\operatorname{Supp}(M)$) est
\emph{équidimensionnel} (16.1.7).

